
Este capítulo descreve a abordagem metodológica adotada para o desenvolvimento do interpretador dinâmico proposto, detalhando a arquitetura geral do sistema, as decisões de projeto que sustentam cada fase do compilador implementado e os mecanismos de integração com a plataforma web. A condução do trabalho seguiu as etapas previstas no cronograma do projeto, contemplando, nesta primeira fase do Trabalho de Conclusão de Curso, a revisão bibliográfica, a especificação do escopo da linguagem e dos comandos dinâmicos, a modelagem da arquitetura do interpretador, a implementação das análises léxica e sintática, a geração e a execução do código intermediário e a integração entre o núcleo do compilador e o ambiente web.

A construção do sistema foi conduzida de forma incremental e iterativa, em consonância com os princípios discutidos por \citeonline{agile2012} e \citeonline{sabbagh2014}. Cada etapa do compilador foi modelada, implementada e validada por meio de testes automatizados antes da agregação à etapa seguinte, em conformidade com a prática de Desenvolvimento Guiado por Testes consolidada por \citeonline{aniche2012} e \citeonline{martin_coder}. Essa abordagem permitiu a verificação contínua de regressões durante a evolução das gramáticas configuráveis e a manutenção da estabilidade das fases de análise mesmo diante das frequentes mudanças impostas pela natureza dinâmica da linguagem proposta.

\section{Modelagem da Arquitetura do Projeto}
\label{sec:modelagem-arquitetura}

A organização do sistema foi concebida sob a ótica da Separação de Responsabilidades defendida por \citeonline{martin_clean}, dividindo o projeto em dois subprojetos independentes e versionados em um único repositório (\textit{monorepo}): o pacote \texttt{compiler}, responsável pelo núcleo do interpretador, e o pacote \texttt{ide}, responsável pelo Ambiente de Desenvolvimento Integrado web e por toda a interação com o usuário. O isolamento entre esses módulos garante que o núcleo de tradução possa evoluir, ser testado e ser executado de maneira autônoma, sem qualquer acoplamento à camada de apresentação, atendendo à premissa de que regras centrais do domínio não devem depender de detalhes periféricos de entrega.

O fluxo de processamento entre os módulos foi estruturado em estágios sequenciais, conforme ilustra a Figura \ref{fig:arquitetura-geral-interpretador}. O código-fonte digitado pelo usuário na IDE é transmitido, juntamente com a configuração corrente da linguagem (palavras-chave personalizadas, delimitadores de bloco, operadores em forma de palavra, literais booleanos e regras de finalização de instrução), ao analisador léxico do pacote \texttt{compiler}. A saída desse analisador, uma lista de \textit{tokens}, alimenta o analisador sintático, que opera por descida recursiva e, de maneira concomitante, aciona o gerador de código intermediário. O conjunto de instruções de três endereços resultante é então recebido pelo interpretador, que executa o programa em tempo real diretamente no servidor da aplicação, repassando as operações de entrada e saída para o terminal embutido na interface.

\begin{figure}[H]
    \centering
    \caption{Fluxo de execução entre os módulos do interpretador dinâmico.}
    \label{fig:arquitetura-geral-interpretador}
    % \includegraphics[width=0.95\textwidth]{Figuras/arquitetura-geral.png}
    \legend{Fonte: elaborado pelo autor.}
\end{figure}

A comunicação entre a IDE e o núcleo do interpretador foi viabilizada por meio de rotas HTTP expostas pela aplicação Next.js. As rotas \texttt{/api/lexer} e \texttt{/api/intermediator} oferecem ao usuário visões intermediárias do processo de tradução, retornando, respectivamente, a tabela de \textit{tokens} reconhecidos e a lista de instruções intermediárias geradas. A rota \texttt{/api/submissions/validate}, por sua vez, executa o ciclo completo de compilação e interpretação para fins de avaliação automática de exercícios, instanciando o analisador léxico, o iterador de \textit{tokens} e o interpretador diretamente no processo do servidor. Essa rota também é responsável por interagir com o serviço de persistência, implementado em FastAPI sobre SQLAlchemy assíncrono, conforme caracterizado por \citeonline{fastapi2026}, mantendo a separação rígida entre a camada de regras de negócio e a infraestrutura de armazenamento.

Para a construção dos componentes de análise, foi adotado o padrão de projeto \textit{Factory} discutido por \citeonline{silveira2012}, materializado pela classe \texttt{LexerScannerFactory}. Esta abstração centraliza a decisão sobre qual analisador especializado deve ser invocado a cada caractere lido, encapsulando a complexidade de instanciação e reduzindo o acoplamento global entre os \textit{scanners} de comentários, identificadores, números, \textit{strings} e símbolos. O mesmo princípio guiou a organização do analisador sintático, no qual cada não-terminal da gramática foi materializado em um arquivo independente, refletindo diretamente a estrutura formal da linguagem e favorecendo tanto a legibilidade quanto a testabilidade do código.

\section{Especificação da Linguagem e Comandos Dinâmicos}
\label{sec:especificacao-linguagem}

A linguagem suportada pelo interpretador foi projetada como um subconjunto didático de linguagens imperativas tipadas, suficientemente expressivo para abranger os conteúdos típicos das disciplinas introdutórias de programação, e simultaneamente flexível para que professores e alunos possam customizar suas regras superficiais. O conjunto base de funcionalidades contempla declarações de variáveis com tipagem estática, atribuições, expressões aritméticas, lógicas e relacionais, estruturas condicionais (\texttt{if}/\texttt{else} e \texttt{switch}), estruturas de repetição (\texttt{for} e \texttt{while}), arranjos uni e multidimensionais com modos de alocação fixos e dinâmicos, e operações de entrada e saída padrão (\texttt{print} e \texttt{scan}).

A definição de quais aspectos da linguagem seriam exibidos como configuráveis foi guiada pelos achados empíricos de \citeonline{stefik2013}, segundo os quais a sintaxe tradicional derivada da família C atua como uma barreira significativa para estudantes iniciantes. Optou-se, portanto, por permitir a personalização precisamente daqueles elementos que, comprovadamente, oneram a carga cognitiva inicial: o vocabulário das palavras reservadas, os operadores lógicos e relacionais (que podem ser substituídos por palavras-chave em linguagem natural), os literais booleanos, os delimitadores de bloco (chaves ou indentação) e o terminador de instrução (ponto e vírgula opcional ou outro caractere arbitrário). Essas opções de configuração foram modeladas pela estrutura \texttt{LexerConfig}, transmitida ao analisador léxico a cada nova compilação e validada antes de qualquer processamento.

A validação prévia da configuração é tratada por funções utilitárias dedicadas, como \texttt{validateOperatorWordMap}, \texttt{validateBooleanLiteralMap}, \texttt{validateBlockDelimiters} e \texttt{validateStatementTerminatorLexeme}. Cada uma dessas funções verifica restrições léxicas elementares antes do início da análise, impedindo, por exemplo, que palavras-chave personalizadas colidam com identificadores válidos ou que delimitadores de bloco contenham caracteres inválidos. Essa estratégia preventiva atende aos princípios de detecção precoce de erros discutidos por \citeonline{cooper2012}, transformando inconsistências de configuração em mensagens claras antes mesmo da leitura do código-fonte.

\section{Interpretador}
\label{sec:metodologia-interpretador}

O núcleo de tradução foi implementado em TypeScript, executado tanto via Node.js no servidor da IDE quanto pelo motor JavaScript do navegador em modos auxiliares de visualização. A escolha pela linguagem TypeScript permitiu o uso intensivo do sistema de tipos durante o desenvolvimento, aproximando o processo de implementação dos benefícios de verificação estática descritos por \citeonline{aho1995} e \citeonline{martin_clean_code}, e contribuindo para a robustez das estruturas de dados que atravessam todas as fases do compilador.

\subsection{Análise Léxica}
\label{subsec:analise-lexica}

A análise léxica foi implementada na classe \texttt{Lexer}, que percorre o código-fonte caractere por caractere e produz uma lista ordenada de instâncias da classe \texttt{Token}. Cada \textit{token} carrega quatro atributos essenciais à fase seguinte e à emissão de erros: um identificador numérico (\texttt{type}), associado a uma das categorias léxicas suportadas pela linguagem, o lexema literalmente lido do código-fonte, a linha e a coluna em que o lexema iniciou-se. Essa estrutura está alinhada à distinção teórica entre lexema e \textit{token} discutida por \citeonline{nystrom2021} e foi mantida estável ao longo de todo o desenvolvimento, pois constitui a interface entre o \textit{front-end} e o restante do compilador.

Os \textit{tokens} foram organizados em sete grupos temáticos --- aritméticos, lógicos, relacionais, atribuições, palavras reservadas, símbolos e literais --- cada um deles definido em um módulo próprio sob o diretório \texttt{token/constants}. Essa segregação favorece a manutenção da tabela de \textit{tokens} e habilita o uso de estilos visuais distintos por categoria na IDE, alinhados aos campos \texttt{TOKENS\_STYLE} associados a cada grupo. A tabela completa de tipos é exportada também na forma invertida (\texttt{BY\_ID}), o que viabiliza a tradução reversa de identificadores numéricos em descrições legíveis, fundamental para fins de depuração e visualização pedagógica.

A escolha por uma análise léxica determinística orientada por caracteres em vez de geradores automáticos como \textit{Lex} foi motivada pela natureza dinâmica das palavras-chave da linguagem proposta. Como o vocabulário não é fixo em tempo de compilação, autômatos finitos pré-gerados perderiam a flexibilidade exigida pelo projeto. Em substituição, adotou-se o padrão \textit{Factory} no diretório \texttt{lexer/scanners}, no qual a função \texttt{LexerScannerFactory.getInstance()} despacha o caractere corrente para o \textit{scanner} apropriado: \texttt{comment.ts} para comentários de linha e de bloco, \texttt{identifier.ts} para identificadores e palavras-chave, \texttt{number.ts} para literais inteiros e em ponto flutuante, \texttt{string.ts} para literais de cadeia com sequências de escape, e \texttt{symbol-and-operator.ts} para os demais símbolos. Um \textit{scanner} adicional, \texttt{statement-terminator.ts}, é instanciado dinamicamente conforme o terminador de instrução configurado pelo usuário, podendo assumir caracteres distintos do ponto e vírgula tradicional ou mesmo ser desativado.

A função \texttt{buildEffectiveKeywordMap}, alimentada pela configuração de personalização proveniente da IDE, constrói em tempo de execução o mapeamento entre os lexemas das palavras-chave correntes e os identificadores numéricos canônicos dos \textit{tokens}. Dessa forma, ao reconhecer um identificador no código-fonte, o analisador léxico verifica primeiro se o lexema corresponde a alguma palavra-chave ativa naquela configuração, classificando-o como \texttt{reserved} quando aplicável e como identificador comum caso contrário. Esse mecanismo concentra toda a flexibilidade linguística no analisador léxico, mantendo os estágios subsequentes inalterados ainda que o vocabulário superficial mude entre execuções.

A análise léxica também é responsável pelo modo de delimitação de blocos. Quando a IDE indica que o usuário optou por blocos baseados em indentação, em substituição às chaves, o \textit{lexer} emite \textit{tokens} sintéticos de \texttt{indent}, \texttt{dedent} e \texttt{newline} em pontos apropriados, aproveitando o registro contínuo de coluna e o tamanho de tabulação configurado pelo usuário. Esse comportamento foi inspirado nas linguagens estruturadas por indentação e permite que a etapa de análise sintática mantenha uma única gramática base, alternando entre blocos delimitados por chaves ou por indentação a partir do conjunto de \textit{tokens} efetivamente gerado.

\subsection{Análise Sintática}
\label{subsec:analise-sintatica}

A análise sintática foi implementada por meio de um analisador descendente recursivo, materializado no diretório \texttt{grammar/syntax}. A escolha por esta estratégia, em detrimento das famílias \textit{bottom-up} discutidas por \citeonline{aho1995}, decorreu de três fatores: a clareza pedagógica da correspondência direta entre regras de produção e funções TypeScript, a facilidade de emissão de mensagens de erro contextualizadas, e a aderência natural ao requisito de visualização passo a passo do processo de compilação dentro da IDE. Cada não-terminal da gramática foi materializado em um arquivo dedicado --- como \texttt{stmt.ts}, \texttt{ifStmt.ts}, \texttt{forStmt.ts}, \texttt{whileStmt.ts}, \texttt{switchStmt.ts}, \texttt{declarationStmt.ts}, \texttt{ioStmt.ts} e \texttt{functionCallExpr.ts} ---, contendo, em comentários, a derivação formal correspondente.

A navegação sobre a sequência de \textit{tokens} produzida pelo analisador léxico foi encapsulada na classe \texttt{TokenIterator}, que oferece operações de \texttt{peek}, \texttt{next}, \texttt{consume} e \texttt{match}, além de mecanismos de consulta ao modo de bloco e ao modo de tipagem corrente. O isolamento dessas operações em um único componente padroniza a forma como cada produção avança sobre o fluxo de \textit{tokens} e protege o restante do analisador de possíveis erros de manipulação direta, como leituras fora dos limites ou consumos indevidos. Esta organização materializa o princípio do Padrão Iterador discutido por \citeonline{silveira2012}.

A gramática base foi cuidadosamente fatorada à esquerda e mantida livre de recursão à esquerda, conforme as restrições necessárias à análise preditiva discutidas por \citeonline{aho1995}. Para acomodar as variações configuráveis pelo usuário, a função \texttt{consumeStmtTerminator} centraliza a verificação do terminador de instrução, podendo aceitar ponto e vírgula, lexemas alternativos definidos pelo usuário ou mesmo nenhum terminador, conforme a configuração ativa. O mesmo princípio governa o tratamento dos delimitadores de bloco: quando a IDE comunica ao analisador o modo indentado, o reconhecimento de \texttt{left\_brace} e \texttt{right\_brace} é substituído pelo casamento dos \textit{tokens} sintéticos \texttt{indent} e \texttt{dedent} gerados pelo \textit{lexer}, mantendo o restante da gramática inalterado.

\subsection{Análise Semântica}
\label{subsec:analise-semantica}

A análise semântica foi entrelaçada com a fase sintática, na forma de tradução dirigida pela sintaxe descrita por \citeonline{cooper2012}. À medida que o analisador descendente reconhece uma construção válida, ele realiza imediatamente as verificações contextuais correspondentes, antes de prosseguir para a próxima produção. Essa estratégia evita a construção explícita e completa de uma árvore sintática intermediária, reduzindo o uso de memória e simplificando a integração com o gerador de código intermediário descrito na Subseção \ref{subsec:geracao-codigo-intermediario}.

Foram contemplados três grupos principais de verificações semânticas. O primeiro refere-se à resolução de nomes e à validação de escopo: cada identificador citado em uma expressão é confrontado com a tabela de símbolos corrente, e referências a variáveis não declaradas são reportadas como erros semânticos. O segundo grupo trata da consistência de tipos em expressões, atribuições, índices de \textit{arrays} e parâmetros de \texttt{scan}; nesse último caso, optou-se pelo registro de um ``hint'' explícito de tipo (\texttt{int} ou \texttt{float}) na instrução de entrada, transferindo a coerção final para o interpretador. O terceiro grupo refere-se à aridade e à estrutura: chamadas de funções com número incompatível de argumentos, acessos a \textit{arrays} com número de dimensões diferente do declarado e usos de \texttt{break} ou \texttt{continue} fora de contextos de repetição são todos sinalizados como erros semânticos contextualizados.

\subsection{Geração de Código Intermediário}
\label{subsec:geracao-codigo-intermediario}

A representação intermediária adotada é uma variante de código de três endereços, conforme caracterizado por \citeonline{aho1995} e \citeonline{cooper2012}, materializada pela interface \texttt{Instruction}, contendo um operador \texttt{op}, um destino \texttt{result} e dois operandos \texttt{operand1} e \texttt{operand2}. O conjunto de operadores suportados é definido como uma união de tipos em \texttt{interpreter/constants.ts} e inclui as seis operações aritméticas (\texttt{+}, \texttt{-}, \texttt{*}, \texttt{/}, \texttt{\%} e \texttt{//}), as três operações lógicas (\texttt{||}, \texttt{\&\&} e \texttt{!}), as seis operações relacionais, atribuições, declarações simples e de \textit{arrays}, instruções de acesso e atualização de \textit{arrays} (\texttt{ARRAY\_GET} e \texttt{ARRAY\_SET}), saltos condicionais e incondicionais (\texttt{IF}, \texttt{JUMP} e \texttt{RETURN}), rótulos (\texttt{LABEL}) e chamadas de sistema (\texttt{CALL}).

A geração dessas instruções é realizada pela classe \texttt{Emitter}, instanciada por cada análise sintática. O \textit{emitter} mantém dois contadores monotônicos responsáveis pela criação de variáveis temporárias (\texttt{\_\_temp0}, \texttt{\_\_temp1}, $\ldots$) e de rótulos (\texttt{\_\_label0}, \texttt{\_\_label1}, $\ldots$), ambos garantidos como únicos no escopo do programa traduzido. As expressões são decompostas em operações binárias elementares, conforme o exemplo clássico apresentado no referencial teórico, e o controle de fluxo das estruturas \texttt{if}, \texttt{while}, \texttt{for} e \texttt{switch} é traduzido em sequências de rótulos e saltos condicionais, eliminando a hierarquia das construções de alto nível em favor da representação linear esperada pelo interpretador.

A opção por gerar três endereços, em detrimento de representações mais abstratas como árvores sintáticas anotadas, foi motivada pelo propósito didático da plataforma. A linearidade das instruções emitidas permite que a IDE apresente, na interface, exatamente o fluxo de operações que será executado, tornando visível para o estudante o vínculo entre comandos de alto nível e suas decomposições elementares. Essa exposição direta da representação intermediária é uma característica diferencial da plataforma em relação às ferramentas analisadas no referencial teórico.

\subsection{Interpretação e Tratamento de Erros em Tempo de Execução}
\label{subsec:interpretacao}

A execução das instruções de três endereços foi implementada na classe \texttt{Interpreter}, projetada para operar sem pilha explícita de execução em sua versão base, com substituição direta de identificadores por seus valores correntes em uma tabela de símbolos. O ponteiro de instruções avança linearmente sobre a lista de \texttt{Instruction}, com exceção dos saltos, que reposicionam o ponteiro para o índice do rótulo associado. Para suportar funções, foi adicionada uma estrutura auxiliar de quadros de chamada (\texttt{CallFrame}), que armazena o endereço de retorno, o nome da variável destino e o escopo local da função invocada, permitindo a execução de chamadas e retornos sem comprometer a simplicidade da máquina virtual.

As operações de entrada e saída foram abstraídas por meio de funções de retorno (\texttt{stdout} e \texttt{stdin}) passadas ao construtor do interpretador. Essa decisão de projeto isola o motor de execução de qualquer detalhe de apresentação, permitindo que, na IDE, as saídas sejam direcionadas ao componente de terminal embutido, enquanto, nos testes automatizados, as mesmas operações sejam capturadas em \textit{buffers} de \textit{string}. Em conformidade com os princípios de baixo acoplamento e alta coesão discutidos por \citeonline{martin_clean}, essa abstração permitiu reutilizar o mesmo interpretador tanto em modo interativo quanto em modo de validação automatizada.

A leitura de valores via \texttt{scan} concentra um dos pontos mais delicados do tratamento de erros em tempo de execução. Como a entrada do usuário pode divergir do tipo esperado pela variável de destino --- por exemplo, um valor textual digitado em uma leitura de inteiro ---, foi implementada uma rotina de coerção, \texttt{coerceValueForType}, complementada pela função \texttt{parseScanInput}, que inspeciona o ``hint'' semântico anotado na instrução e tenta converter a cadeia lida para o tipo declarado. Quando a conversão é impossível, em vez de propagar uma exceção genérica de \textit{runtime}, o interpretador lança uma instância de \texttt{RuntimeError} contendo a posição original do comando no código-fonte e uma mensagem traduzida via o módulo \texttt{i18n}, permitindo que a IDE destaque visualmente a linha responsável e exiba o erro no idioma corrente do usuário.

Os erros estáticos detectados durante a compilação foram modelados pelo módulo \texttt{issue}, com três categorias distintas: \texttt{IssueError} representa falhas que interrompem a tradução, \texttt{IssueWarning} representa alertas não fatais coletados durante a análise, e \texttt{IssueInfo} acumula mensagens informativas relevantes ao usuário. Toda ocorrência registra obrigatoriamente um código identificador, uma mensagem, a linha e a coluna do problema, permitindo que a IDE marque o trecho exato no editor e ofereça orientações pedagógicas associadas ao código do erro.

\section{Integração entre a Plataforma e o Interpretador}
\label{sec:integracao-plataforma}

A integração entre a plataforma web e o núcleo do interpretador foi materializada por meio do consumo direto do pacote \texttt{compiler} como dependência local do pacote \texttt{ide}, sem qualquer empacotamento intermediário em formato binário ou em servidor remoto separado. As importações realizadas na rota \texttt{/api/submissions/validate} acessam diretamente as classes \texttt{Lexer}, \texttt{TokenIterator} e \texttt{Interpreter}, bem como os tipos de configuração léxica, garantindo que qualquer alteração no compilador se propague à IDE sem etapas adicionais de implantação.

O ciclo de submissão de exercícios foi modelado em quatro passos sequenciais. Primeiro, a configuração corrente da linguagem é normalizada pela função \texttt{normalizeCompilerConfig}, que reconcilia o estado salvo do exercício com a configuração ativa na IDE e produz um descritor de gramática \texttt{IDEGrammarConfig} consistente. Em seguida, para cada caso de teste, instancia-se um analisador léxico configurado com esse descritor e executa-se a tradução completa do código submetido. Posteriormente, o interpretador é instanciado com funções \texttt{stdout} e \texttt{stdin} ligadas a \textit{buffers} controlados, permitindo que a saída produzida seja comparada com a saída esperada e que a entrada simulada do caso de teste seja consumida sequencialmente. Por fim, o resultado consolidado de cada caso é encapsulado em uma estrutura \texttt{TTestCaseResult} e agregado em um \texttt{TValidationResult}, devolvido ao \textit{frontend} para apresentação ao usuário.

A persistência de exercícios, turmas e submissões foi mantida em um serviço externo, implementado em Python sobre o \textit{framework} FastAPI e o ORM SQLAlchemy assíncrono, com versionamento de \textit{schema} via Alembic, em conformidade com os princípios de abstração de persistência discutidos por \citeonline{martin_clean}. Essa separação garante que o núcleo de compilação e interpretação permaneça agnóstico em relação à camada de armazenamento, podendo ser executado, inclusive, sem qualquer dependência de banco de dados nas execuções locais e nos testes automatizados.

\section{Estratégia de Testes Automatizados}
\label{sec:metodologia-testes}

A validação contínua das fases do compilador foi sustentada por uma suíte de testes automatizados construída com o \textit{framework} Vitest, organizada no diretório \texttt{packages/compiler } \texttt{/src/tests} e estruturada de modo a refletir a topologia interna do projeto. Os testes foram subdivididos em três camadas: testes de \textit{tokens}, que verificam a estabilidade dos identificadores numéricos e dos agrupamentos por categoria; testes de \textit{lexer}, que exercitam cada \textit{scanner} especializado em isolamento, abrangendo comentários, literais numéricos, cadeias com escapes, terminadores configuráveis, delimitadores de bloco, operadores em forma de palavra e geração de \textit{tokens} de indentação; e testes de gramática, que validam a aceitação ou rejeição de programas inteiros sob distintas configurações de tipagem, modos de bloco e variantes do comando \texttt{switch}.

A adoção do TDD discutida por \citeonline{aniche2012} mostrou-se particularmente útil em construções como o terminador de instrução configurável e o modo de blocos por indentação, em que pequenas alterações na lógica do \textit{lexer} poderiam quebrar inadvertidamente programas que antes compilavam corretamente. Em conformidade com a pirâmide de testes discutida por \citeonline{martin_clean}, a maior parte da suíte concentra-se no nível unitário, garantindo execuções rápidas e \textit{feedback} imediato, enquanto os testes de gramática operam como testes de integração internos ao compilador, verificando o comportamento agregado das fases léxica, sintática e de geração de código intermediário.

Os mesmos testes automatizados são executados localmente durante o desenvolvimento e também acionados pelo \textit{pipeline} de integração contínua antes da incorporação de qualquer mudança ao ramo principal do repositório, assegurando que as características essenciais do compilador --- estabilidade dos identificadores de \textit{tokens}, correção semântica das instruções intermediárias geradas e equivalência de comportamento entre as variantes configuráveis da linguagem --- sejam preservadas ao longo de toda a evolução incremental do projeto.
